\documentclass[10pt]{article}

\usepackage{parskip}
\usepackage[margin=1in]{geometry} 
\usepackage{amsmath,amsthm,amssymb, graphicx, multicol, array}
\usepackage{enumitem}
\usepackage{amssymb}
\usepackage{float}
\usepackage{href-ul}
 
\newcommand{\N}{\mathbb{N}}
\newcommand{\Z}{\mathbb{Z}}
 
\newenvironment{problem}[2][Problem]{\begin{trivlist}
\item[\hskip \labelsep {\bfseries #1}\hskip \labelsep {\bfseries #2.}]}{\end{trivlist}}

\begin{document}
 
\title{Mixed Integer Programming}
\author{Jakob Sverre Alexandersen\\
GRA4154 Supply Chain Optimization with Mathematical Programming\\
Lecture 2}
\maketitle

\tableofcontents
\newpage

\section{Binary variables}

\begin{itemize}
    \item Can take the values 0 and 1 only
    \item For most real-world applications, binary variables are necessary in order to obtain a realistic model 
\end{itemize}

\textbf{Example: an investment problem}
\begin{itemize}
    \item A company can choose between the following investment projects: 1, …, 8
    \item In total, the company can invest up to 100M
    \item The company wants to maximize total NPV
    \item Each project is og the type ``all in or nothing'', i.e., partial investment in the project is not possible
    \item This is a variation of the knapsack problem
\end{itemize}

\subsection{Problem with only binary variables}
A manufacturer of electronic components has limited production capacity and needs to prioritize which customer orders to satisfy during periods of high demand. 

The company wants to select orders so as to maximize total profit. Each order must either be satisfied completely or rejected. 

On a given day, the manufacturer must choose between the following orders: 
\begin{center}
    \begin{tabular}{|c|c|c|c|}
        \hline 
        & \textbf{Revenue} & \textbf{Cost} & \textbf{Hours per order}\\\hline 
        \textbf{Customer 1} & 18 & 8 & 3 \\\hline
        \textbf{Customer 2} & 32 & 9 & 7 \\\hline
        \textbf{Customer 3} & 25 & 9 & 5 \\\hline
        \textbf{Customer 4} & 24 & 10 & 4 \\\hline
        \textbf{Customer 5} & 26 & 13 & 3 \\\hline
    \end{tabular}
\end{center}

``Hours per unit'' is the number of production hours spent for the order. 

There are 16 hours available in the production facility

\textbf{Parameters:}
\begin{itemize}
    \item $r_i $: revenue per order 
    \item $c_i $: cost per order 
    \item $a_i $: hours per order 
    \item $b $: capacity
\end{itemize}

\textbf{Variables}
\[
X_i = 
\begin{cases}
    0 \quad\text{if we reject order }i\\
    1\quad\text{otherwise}
\end{cases}
\]
\textbf{Model}
\begin{gather*}
    \max\left(\sum_i (r_i - c_i)X_i\right)\\
    \text{s.t.}\\
    \sum_i a_iX_i \le b \\
    X_i \in \{0, 1\}
\end{gather*}
\newpage 
\subsection{Model with both continuous and binary variables}
\textbf{Mixed-Integer Problem: MIP}

A company needs to decide production quantities of five product types so to maximize total profit. 

For each product type, there is a market demand that defines the maximum possible quantity sold 

\textbf{Now we have added} a fixed cost for each product type. The cost is avoided if the product type is not produced

The five product types share a joint production capacity that limits the total production:

\begin{center}
    \begin{tabular}{|c|c|c|c|c|c|}
        \hline 
        & \textbf{Sales price } & \textbf{Variable cost } & \textbf{Hours per } & \textbf{Max number of } & \textbf{Fixed cost}\\
        & \textbf{per unit} & \textbf{per unit} & \textbf{unit produced} & \textbf{units sold} & \\\hline
        \textbf{Type 1} & 18 & 8 & 0.30 & 2500  & 3700 \\\hline
        \textbf{Type 2} & 32 & 9 & 0.70 & 3000  & 1900\\\hline
        \textbf{Type 3} & 25 & 9 & 0.50 & 4000  & 1000\\\hline
        \textbf{Type 4} & 24 & 10 & 0.40 & 2000 & 2200 \\\hline
        \textbf{Type 5} & 26 & 13 & 0.30 & 1500 & 2600 \\\hline
    \end{tabular}
\end{center}
There are 2000 hours available for production 

\textbf{Parameters}
\begin{itemize}
    \item $p_i $: price per unit 
    \item $c_i $: cost per unit 
    \item $h_i $: hours per unit
    \item $m_i $: max amount per unit
    \item $b $: capacity 
    \item $\mathcal{F}_i$: fixed cost
\end{itemize}

\textbf{Variables}
\begin{itemize}
    \item $X_i$: quantity
\end{itemize}

\[
\lambda_i = 
\begin{cases}
    0\quad\text{if type is not produced}\\
    1\quad\text{otherwise}
\end{cases}
\]

\textbf{Model}
\begin{gather*}
    \max\left(\sum_i \big((p_i - c_i)X_i + \mathcal{F}_i\big)\lambda_i\right)\\
    \text{s.t.}\\
    0 \le X_i \le m_i\quad\forall i \\
    \sum_i h_iX_i \le b\quad\forall i\\
    X_i + \lambda_i
\end{gather*}

\end{document}